\section*{NeurIPS Paper Checklist}

\begin{enumerate}

\item {\bf Claims}
    \item[] Question: Do the main claims made in the abstract and introduction accurately reflect the paper's contributions and scope?
    \item[] Answer: \answerYes{}
    \item[] Justification: The abstract and introduction state the headline empirical claim (hand-crafted features outperform every evaluated TSFM on anomaly probing, while four other properties saturate non-model baselines under the canonical protocol), and the methodological claim (representation-level claims about TSFMs require explicit non-model controls). Both are directly supported by the canonical results table in the main body and by the appendix manifest, retroactive case study, adversarial probing, and seed-bootstrap LEACE artifacts.

\item {\bf Limitations}
    \item[] Question: Does the paper discuss the limitations of the work performed by the authors?
    \item[] Answer: \answerYes{}
    \item[] Justification: The paper discusses limitations including PCA-compressed views for some models, protocol-level rather than raw-state cross-family comparisons, and weak synthetic-to-real transfer; Appendix A.1 also states scope limitations explicitly.

\item {\bf Theory assumptions and proofs}
    \item[] Question: For each theoretical result, does the paper provide the full set of assumptions and a complete (and correct) proof?
    \item[] Answer: \answerNA{}
    \item[] Justification: This submission is an empirical benchmark paper and does not claim new theorems or proofs.

\item {\bf Experimental result reproducibility}
    \item[] Question: Does the paper fully disclose all the information needed to reproduce the main experimental results of the paper to the extent that it affects the main claims and/or conclusions of the paper (regardless of whether the code and data are provided or not)?
    \item[] Answer: \answerYes{}
    \item[] Justification: Primary confirmatory results use a canonical 60/20/20 train/val/test protocol with best-layer selection on val, final reporting on held-out test, matched sklearn estimators, 5 seeds, and bootstrap 95\% CIs. Known preprocessing caveats (PCA fit scope, extraction heterogeneity across model families) are disclosed in the limitations and appendix. Legacy PyTorch probing results (80/10/10 split) are retained only as auxiliary diagnostics.

\item {\bf Open access to data and code}
    \item[] Question: Does the paper provide open access to the data and code, with sufficient instructions to faithfully reproduce the main experimental results, as described in supplemental material?
    \item[] Answer: \answerYes{}
    \item[] Justification: An anonymized code artifact is provided as supplementary material, containing all source code, test suite (170 tests), synthetic data generators, extraction pipelines, and evaluation scripts with reproduction instructions.

\item {\bf Experimental setting/details}
    \item[] Question: Does the paper specify all the training and test details (e.g., data splits, hyperparameters, how they were chosen, type of optimizer) necessary to understand the results?
    \item[] Answer: \answerYes{}
    \item[] Justification: The paper reports the canonical confirmatory setup (60/20/20 train/val/test, matched sklearn LogisticRegression/Ridge estimators, 5 seeds, bootstrap 95\% CIs), the temporal properties and datasets, and the intervention settings. Architecture details and supplementary settings are in the appendix. Baseline controls and model probes share the identical scoring pipeline in the canonical protocol; legacy PyTorch probing at 80/10/10 is retained only for auxiliary diagnostics and is disclosed as such.

\item {\bf Experiment statistical significance}
    \item[] Question: Does the paper report error bars suitably and correctly defined or other appropriate information about the statistical significance of the experiments?
    \item[] Answer: \answerYes{}
    \item[] Justification: Primary confirmatory results (canonical baseline-control table and per-model canonical rerun) report 5-seed bootstrap 95\% confidence intervals on held-out test scores; seed-bootstrap LEACE results for the four canonical encoder-side models are also reported in the appendix. The legacy 3-seed PyTorch probe table is retained only as auxiliary diagnostic and is labeled as such. Cross-property entanglement, downstream impact, and other single-run results are explicitly labeled exploratory and are not used as primary evidence. Real-world comparisons with small $n$ are reported as reference diagnostics rather than rankings.

\item {\bf Experiments compute resources}
    \item[] Question: For each experiment, does the paper provide sufficient information on the computer resources (type of compute workers, memory, time of execution) needed to reproduce the experiments?
    \item[] Answer: \answerYes{}
    \item[] Justification: The appendix states that the benchmark was run on a single NVIDIA A100 80GB GPU and reports the approximate end-to-end compute budget for extraction, probing, intervention, and similarity analysis.

\item {\bf Code of ethics}
    \item[] Question: Does the research conducted in the paper conform, in every respect, with the NeurIPS Code of Ethics \url{https://neurips.cc/public/EthicsGuidelines}?
    \item[] Answer: \answerYes{}
    \item[] Justification: The work analyzes public models and public benchmark datasets only,
    does not involve private or human-subject data, and explicitly frames intervention methods as
    diagnostic tools rather than undisclosed deployment-time manipulation mechanisms.

\item {\bf Broader impacts}
    \item[] Question: Does the paper discuss both potential positive societal impacts and negative societal impacts of the work performed?
    \item[] Answer: \answerYes{}
    \item[] Justification: The paper discusses interpretability benefits for safer model diagnosis and selection, while also noting that steering/erasure techniques could be misused if applied outside controlled evaluation settings.

\item {\bf Safeguards}
    \item[] Question: Does the paper describe safeguards that have been put in place for responsible release of data or models that have a high risk for misuse (e.g., pre-trained language models, image generators, or scraped datasets)?
    \item[] Answer: \answerNA{}
    \item[] Justification: The paper does not release new generative models or scraped datasets; it evaluates already public time-series models and public benchmark data.

\item {\bf Licenses for existing assets}
    \item[] Question: Are the creators or original owners of assets (e.g., code, data, models), used in the paper, properly credited and are the license and terms of use explicitly mentioned and properly respected?
    \item[] Answer: \answerYes{}
    \item[] Justification: The paper cites the original model and dataset papers and now states in
    the appendix that the benchmark uses only public datasets and checkpoints under their original
    access terms, with no deprecated or private datasets introduced by this submission.

\item {\bf New assets}
    \item[] Question: Are new assets introduced in the paper well documented and is the documentation provided alongside the assets?
    \item[] Answer: \answerYes{}
    \item[] Justification: An anonymized code artifact is provided as supplementary material with a README containing installation and reproduction instructions, 170 automated tests, and a Benchmark Datasheet (\Cref{sec:datasheet}) following the Datasheets for Datasets framework. The benchmark includes synthetic data generators, model wrappers, and evaluation scripts.

\item {\bf Crowdsourcing and research with human subjects}
    \item[] Question: For crowdsourcing experiments and research with human subjects, does the paper include the full text of instructions given to participants and screenshots, if applicable, as well as details about compensation (if any)?
    \item[] Answer: \answerNA{}
    \item[] Justification: The work does not involve crowdsourcing or human subjects.

\item {\bf Institutional review board (IRB) approvals or equivalent for research with human subjects}
    \item[] Question: Does the paper describe potential risks incurred by study participants, whether such risks were disclosed to the subjects, and whether Institutional Review Board (IRB) approvals (or an equivalent approval/review based on the requirements of your country or institution) were obtained?
    \item[] Answer: \answerNA{}
    \item[] Justification: The work does not involve human subjects.

\item {\bf Declaration of LLM usage}
    \item[] Question: Does the paper describe the usage of LLMs if it is an important, original, or non-standard component of the core methods in this research? Note that if the LLM is used only for writing, editing, or formatting purposes and does \emph{not} impact the core methodology, scientific rigor, or originality of the research, declaration is not required.
    \item[] Answer: \answerNA{}
    \item[] Justification: LLMs are not part of the benchmark methodology, model training, probing, intervention, or analysis pipeline described in the paper.

\end{enumerate}
